\subsection{From sequence limits to function limits}

A sequence approaches its limit through its tail. For a function, the input may approach a point through infinitely many different values, so we first state the local definition that every route must satisfy.

\begin{definitionbox}[Function limit]
Let $f\colon D\to\mathbb R$ and let $a$ be an accumulation point of $D$. We write
\[
\lim_{x\to a}f(x)=L
\]
when, for every $\varepsilon>0$, there exists $\delta>0$ such that for every $x\in D$,
\[
0<|x-a|<\delta
\quad\Longrightarrow\quad
|f(x)-L|<\varepsilon.
\]
\end{definitionbox}

The condition excludes $x=a$: a function limit is determined by nearby inputs, not by the value assigned at the point itself.

A function limit can also be tested by following the function along every sequence of inputs approaching the point. This gives the natural bridge from the discrete theory of sequences to local function behaviour.

\begin{theorembox}[Heine's theorem]
Let $a$ be an accumulation point of the domain of $f$. Then
\[
\lim_{x\to a}f(x)=L
\]
if and only if, for every sequence $(x_n)$ in the domain of $f$ satisfying
\[
x_n\ne a
\qquad\text{and}\qquad
x_n\to a,
\]
we have
\[
f(x_n)\to L.
\]
\end{theorembox}

\begin{remarkbox}[Proof status]
The forward implication follows by applying the $\varepsilon$--$\delta$ definition to sufficiently late terms $x_n$. The reverse implication is proved by contradiction: if the local definition failed, one could choose inputs successively closer to $a$ whose outputs remain a fixed distance from $L$. Those inputs would form a sequence contradicting the sequential condition. The full quantified proof is not expanded here.
\end{remarkbox}

\begin{examplebox}[Using Heine's theorem]
Consider
\[
f(x)=2x+1.
\]
Take any sequence $x_n\to3$. Then
\[
f(x_n)=2x_n+1\to2\cdot3+1=7.
\]
Because the conclusion holds for every sequence approaching $3$,
\[
\lim_{x\to3}(2x+1)=7.
\]
\end{examplebox}

\begin{examplebox}[How to disprove a function limit]
For
\[
f(x)=\frac{x}{|x|},\qquad x\ne0,
\]
choose
\[
x_n=\frac1n
\qquad\text{and}\qquad
y_n=-\frac1n.
\]
Both sequences approach $0$, but
\[
f(x_n)=1,
\qquad
f(y_n)=-1.
\]
The output sequences have different limits, so $\lim_{x\to0}f(x)$ does not exist.
\end{examplebox}

\begin{interpretationbox}[A function limit must survive every route]
The local definition controls all sufficiently near inputs at once. Heine's theorem expresses the same requirement through every admissible approaching sequence. One chosen sequence can disprove a limit when it behaves differently, but it cannot prove a limit by itself.
\end{interpretationbox}
